\documentclass[11pt,a4paper]{article}

\usepackage[utf8]{inputenc}
\usepackage[T1]{fontenc}
\usepackage[french]{babel}
\usepackage[a4paper,margin=1.8cm]{geometry}
\usepackage{amsmath,amssymb}
\usepackage{lmodern}
\usepackage{fancyhdr}
\usepackage{enumitem}

\pagestyle{fancy}
\fancyhf{}
\fancyhead[L]{3\ieme}
\fancyhead[C]{Entraînement Brevet – Nombres et Calculs}
\fancyhead[R]{\thepage}

\title{\Huge 40 Exercices contextualisés + 4 DS \\ Niveau Brevet}
\date{}

\begin{document}
\maketitle

%====================================================
\section*{PARTIE 1 – EXERCICES CONTEXTUALISÉS}

\subsection*{★ Niveau 1}

\begin{enumerate}

\item Une pizza est partagée en 8 parts. Paul en mange $\frac{3}{8}$ et Léa $\frac{2}{8}$.  
Quelle fraction de pizza a été mangée ?

\item Un magasin applique une réduction de 10\% sur un article à 50€.  
Quel est le prix final ?

\item Une voiture consomme 5 L pour 100 km.  
Combien consomme-t-elle pour 200 km ?

\item Un carré a une aire de 36 m².  
Quelle est la longueur de son côté ?

\item Un club compte 20 élèves dont $\frac{3}{4}$ sont des filles.  
Combien y a-t-il de filles ?

\end{enumerate}

%====================================================
\subsection*{★★ Niveau 2}

\begin{enumerate}[resume]

\item Un cinéma propose un tarif : 8€ + 2€ par film.  
Exprimer le prix en fonction du nombre de films.

\item Un terrain rectangulaire mesure 12 m sur 5 m.  
Calculer son aire.

\item Un téléphone coûte 300€ avec une réduction de 20\%.  
Calculer le prix final.

\item Une classe a 30 élèves dont 60\% sont des garçons.  
Combien de garçons ?

\item Simplifier une recette : 6 œufs pour 4 personnes.  
Combien d’œufs pour 2 personnes ?

\end{enumerate}

%====================================================
\subsection*{★★★ Niveau 3}

\begin{enumerate}[resume]

\item Une entreprise facture 15€ + 3€ par heure.  
Combien coûte 10h ?  
Pour combien d’heures paie-t-on 60€ ?

\item Un sac contient 10 billes rouges et 5 bleues.  
Quelle est la probabilité de tirer une rouge ?

\item Décomposer 84 en facteurs premiers et expliquer à quoi cela sert dans un partage.

\item Résoudre :  
Un nombre augmenté de 5 puis multiplié par 2 donne 30.

\item Une aire est donnée par $(x+2)(x+3)$.  
Développer et interpréter.

\end{enumerate}

%====================================================
\subsection*{★★★★ Niveau 4}

\begin{enumerate}[resume]

\item Une entreprise vend des objets : coût = $2x + 50$.  
Calculer pour 20 objets et résoudre pour 150€.

\item Un escalier a des marches de 17 cm.  
Pour atteindre 272 cm, combien de marches ?

\item On partage 84 bonbons rouges et 126 bleus.  
Former des sachets identiques : combien ?

\item Simplifier une expression :  
$(2x+3)^2 - (x+1)(x+5)$

\item Une population double chaque jour.  
Combien après 10 jours ?

\end{enumerate}

%====================================================
\subsection*{★★★★★ Niveau 5}

\begin{enumerate}[resume]

\item Deux programmes de calcul donnent :  
$-2x+5$ et $3x-4$.  
Trouver quand ils donnent le même résultat.

\item Une entreprise compare deux tarifs :  
A : 10€ + 2x  
B : 5€ + 3x  
À partir de combien d’unités B devient plus cher ?

\item Montrer que $(n+1)^2 - n^2 = 2n+1$ et interpréter.

\item Décomposer 360 et expliquer comment répartir équitablement.

\item Simplifier :
\[
\frac{3}{4} + \frac{5}{6} \times \frac{8}{3}
\]

\end{enumerate}

%====================================================
%====================================================
\newpage
\section*{DS N°1 – Type Brevet (100 points)}

\textbf{Exercice 1 – Calcul numérique (20 pts)}  
\[
A = \frac{3}{5} + \frac{7}{10} \times \frac{4}{3}
\]

\textbf{Exercice 2 – Calcul littéral (20 pts)}  
Développer et réduire :
\[
(2x+3)^2 - (x+1)(x+5)
\]

\textbf{Exercice 3 – Équation (20 pts)}  
Un abonnement coûte 10€ + 2€/h.  
1. Exprimer le prix  
2. Calculer pour 8h  
3. Résoudre pour 30€

\textbf{Exercice 4 – Arithmétique (20 pts)}  
On a 84 bonbons rouges et 126 bleus.  
Former des sachets identiques.

\textbf{Exercice 5 – Programme de calcul (20 pts)}  
Programme :  
Multiplier par 2 → ajouter 5 → multiplier par 3  
Exprimer en fonction de $x$

%====================================================
\newpage
\section*{DS N°2 – Téléphonie et croissance (100 points)}

\textbf{Exercice 1 – Forfaits téléphoniques (20 pts)}  
Deux offres :  
A : 10€ + 2x  
B : 5€ + 3x  

1. Calculer pour 5 appels  
2. Résoudre A = B  
3. Interpréter  

\bigskip

\textbf{Exercice 2 – Croissance (20 pts)}  
Une population double chaque jour.  
1. Exprimer le nombre après $n$ jours  
2. Calculer après 8 jours  
3. Comparer avec une croissance linéaire

\bigskip

\textbf{Exercice 3 – Calcul littéral (20 pts)}  
\[
B = (x+2)(x+5) - (x+1)^2
\]

1. Développer  
2. Réduire  

\bigskip

\textbf{Exercice 4 – Arithmétique (20 pts)}  
Décomposer 180 et 84.  
Calculer leur PGCD.  
Interpréter dans un problème de partage.

\bigskip

\textbf{Exercice 5 – Équation (20 pts)}  
Résoudre :
\[
3x + 5 = 2x + 12
\]

%====================================================
\newpage
\section*{DS N°3 – Entreprise et modélisation (100 points)}

\textbf{Exercice 1 – Chiffre d’affaires (20 pts)}  
Une entreprise vend des objets :  
Prix = $2x + 50$

1. Calculer pour 20 objets  
2. Résoudre pour 150€  
3. Interpréter  

\bigskip

\textbf{Exercice 2 – Calcul numérique (20 pts)}  
\[
A = \frac{5}{6} - \frac{2}{3} + \frac{3}{4}
\]

\bigskip

\textbf{Exercice 3 – Identité remarquable (20 pts)}  
Développer :
\[
(x+3)^2
\]

\bigskip

\textbf{Exercice 4 – Inéquation (20 pts)}  
Résoudre :
\[
2x + 3 > 7
\]

\bigskip

\textbf{Exercice 5 – Programme de calcul (20 pts)}  
Programme :  
Choisir un nombre  
Soustraire 5  
Multiplier par 3  
Ajouter 11  

1. Calculer pour $x=2$  
2. Exprimer en fonction de $x$  
3. Résoudre pour 20  

%====================================================
\newpage
\section*{DS N°4 – Niveau difficile (100 points)}

\textbf{Exercice 1 – Calcul complexe (20 pts)}  
\[
A = \frac{3}{4} + \frac{5}{6} \times \frac{8}{3}
\]

\bigskip

\textbf{Exercice 2 – Démonstration (20 pts)}  
Montrer que :
\[
(n+1)^2 - n^2 = 2n+1
\]

\bigskip

\textbf{Exercice 3 – Équation (20 pts)}  
Résoudre :
\[
-2x + 5 = 3x - 10
\]

\bigskip

\textbf{Exercice 4 – Arithmétique avancée (20 pts)}  
Décomposer 360 et 252.  
Calculer le PGCD.

\bigskip

\textbf{Exercice 5 – Problème ouvert (20 pts)}  
Deux tarifs :  
A : 10€ + 2x  
B : 5€ + 3x  

1. Comparer  
2. À partir de combien B devient plus cher  
3. Interpréter  

\end{document}